\documentclass[11pt]{article}
\usepackage[T1]{fontenc}
\usepackage[a4paper,margin=1in]{geometry}
\usepackage{amsmath,amssymb,amsthm}
\usepackage[hidelinks]{hyperref}
\usepackage{microtype}

\newtheorem*{sourceproblem}{Source program}
\newtheorem*{laterresult}{Later partial theorem}

\setlength{\parindent}{0pt}
\setlength{\parskip}{0.58em}

\title{Literature Status: A Quantitative Local Advance on\\
Planar $W^{2,1}$ Diffeomorphic Approximation}
\author{Run \texttt{fa\_banach\_001} --- lane 1}
\date{13 August 2026}

\begin{document}
\maketitle

\begin{center}
\fbox{\parbox{0.9\textwidth}{
\textbf{Status: literature-implied answer to a partial subcase; the full
problem remains open.}
Campbell--Hencl ask for an extension step that would turn their approximation
outside a small exceptional set into a global $W^{2,1}$ diffeomorphic
approximation.  D'Onofrio, arXiv:2604.04592, proves a quantitative localized
smoothing theorem for a substantially more structured piecewise-quadratic
class, but explicitly leaves both a finite-singularity completion step and
the general geometric approximation step unresolved.
}}
\end{center}

\section{The source program}

D. Campbell and S. Hencl, \emph{Approximation of planar Sobolev
$W^{2,1}$ homeomorphisms by piecewise quadratic homeomorphisms and
diffeomorphisms}, arXiv:2008.05736, prove that a planar
$W^{2,1}$ homeomorphism with positive Jacobian almost everywhere can be
approximated by a smooth diffeomorphism outside a locally finite union of
squares of arbitrarily small area.  On PDF page~4 they state their next
step~\cite{source}:

\begin{sourceproblem}
Obtain a second-order analogue of the key planar extension theorem used in
the $W^{1,1}$ theory, thereby filling the exceptional squares and obtaining
the full $W^{2,1}$ approximation theorem.  A further goal is to remove the
assumption $J_f>0$ almost everywhere.
\end{sourceproblem}

The first clause is the target considered here.

\section{The 2026 partial advance}

L. D'Onofrio, \emph{$W^{2,1}$ approximation of planar Sobolev
homeomorphisms by smooth diffeomorphisms}, arXiv:2604.04592, begins by
stating that no general $W^{2,1}$ approximation theorem is known.  Its
Assumption~5.1 (PDF page~14) concerns a bounded polygonal domain with a
finite conforming rectangular partition and a map $g$ which is

\begin{itemize}
\item quadratic on each cell and $C^1$-compatible across every interface;
\item a homeomorphism onto its image;
\item quantitatively nondegenerate, with $\det Dg\geq\lambda>0$; and
\item quantitatively injective, with
      $|g(x)-g(y)|\geq m|x-y|$ for all $x,y$.
\end{itemize}

Theorem~5.5 (PDF pages~17--20) then gives the following~\cite{later}.

\begin{laterresult}
For every $\delta>0$, there is a $C^1$ map
$\widetilde g_\delta$ with
\[
 \|\widetilde g_\delta-g\|_{W^{2,1}}<\delta,
 \qquad
 \|D\widetilde g_\delta-Dg\|_{L^\infty}<\delta,
 \qquad
 \det D\widetilde g_\delta\geq\lambda/2.
\]
Moreover,
\[
 |\widetilde g_\delta(x)-\widetilde g_\delta(y)|
 \geq (m-Q_\Omega\delta)|x-y|,
\]
so it is injective for sufficiently small $\delta$, and it is smooth away
from arbitrarily small neighborhoods of the finite set of endpoints of the
interior edges.
\end{laterresult}

This resolves a genuine local analytic part of the Campbell--Hencl program:
interface and vertex regularization can preserve $W^{2,1}$ control,
positive Jacobian, and global injectivity under explicit quantitative
hypotheses.

\section{Why this is not the full answer}

There are two strict gaps between the later theorem and the source problem.
First, the source's approximation outside small squares does not furnish a
global finite, $C^1$-compatible, uniformly bi-Lipschitz piecewise-quadratic
map satisfying Assumption~5.1.  The 2026 paper itself calls construction of
such structured approximants the remaining ``global geometric
approximation problem.''

Second, Theorem~5.5 is smooth only off small neighborhoods of finitely many
edge endpoints.  Its Theorem~6.2 obtains globally smooth injective maps only
after imposing the separate, unproved Assumption~6.1, which is precisely a
localized completion assertion.  Thus arXiv:2604.04592 is a substantial
partial advance and an explicit reduction, not a resolution of the general
question.  No new mathematical result is claimed in this packet.

\begin{thebibliography}{9}
\bibitem{source}
D. Campbell and S. Hencl,
\emph{Approximation of planar Sobolev $W^{2,1}$ homeomorphisms by
piecewise quadratic homeomorphisms and diffeomorphisms},
ESAIM Control Optim. Calc. Var. \textbf{27} (2021), article 26;
arXiv:2008.05736.

\bibitem{later}
L. D'Onofrio,
\emph{$W^{2,1}$ approximation of planar Sobolev homeomorphisms by smooth
diffeomorphisms}, arXiv:2604.04592 (2026).
\end{thebibliography}

\end{document}
